% Tutorial example: data-enabled predictive control: ex9_DeePC/9.9_summary_figures.ipynb (figures), 9.7--9.8 (supporting studies)

\begin{tutorialexample}[label={ex:deepc}]
    {From data to control without a model}
    {TODO-notebook-url}
  Willems' fundamental lemma turns one recorded trajectory into a predictor, so the questions that decide whether DeePC works in practice are how much data the Hankel matrix needs and what happens when the data are not generated by an \ac{LTI} system. This example answers both on the mountain car and shows that regularization---not more data---is what carries the method beyond the ideal case. The complete implementation is provided in \cite[Ex.~9.7--9.9]{cfr-nootbook}.

  We control the car through its position only, $y=p$ ($n_\mathrm{u}=1$, $n_\mathrm{x}=2$, lag $l=2$), sampled at $20\,\mathrm{Hz}$, and solve~\eqref{DDPC_problem} with horizon $L=20$, an initialization window of $T_\mathrm{ini}=4\ge l$ samples in~\eqref{initial_condition_DDPC}, the cost $\ell(u,y)=(y-p^\mathrm{s})^2+0.1\,u^2$, and $|u|\le 1$. The offline input is a uniformly random excitation of length~$N$. Performance is measured by the RMS deviation of the closed-loop position from that of an \ac{MPC} that uses the exact model of the respective plant, so a value of zero means DeePC and model-based control coincide. On flat terrain the plant is \ac{LTI}, and Theorem~\ref{thmdtwill} requires the input to be \ac{PE} of order $L+T_\mathrm{ini}+n_\mathrm{x}$, which by counting columns of the Hankel matrix gives the data-length bound $N\ge(n_\mathrm{u}+1)(L+T_\mathrm{ini}+n_\mathrm{x})-1=51$. \autoref{fig:tutorial-deepc-flat} confirms that this bound is sharp in practice: for shorter records the columns of~\eqref{prediction_model_DDPC} cannot span the measured history, every QP is infeasible and the car does not move; from $N=51$ on\footnote{The \ac{PE} condition is sufficient, not necessary. For the generic random input used here the rank condition~\eqref{generalcond} already holds two samples earlier, at $N=49$, where the Hankel matrix has exactly $n_\mathrm{u}(L+T_\mathrm{ini})+n_\mathrm{x}$ columns \cite{MarkovskyDor2023}.} the closed loop reproduces the exact-model \ac{MPC} to solver precision ($10^{-8}\,\mathrm{m}$), and longer records change nothing but the size of the QP.

  On bumpy terrain, $h(p)=k\cos(18p)$, the plant is nonlinear and the data recorded on it are not consistent with any \ac{LTI} behaviour: the Hankel matrix becomes full rank exactly as it would under measurement noise, the equality~\eqref{prediction_model_DDPC} can match any history, and the unregularized coefficient~$\alpha$ carries no information about the plant. We therefore use the regularized form, replacing $\bar y_{[-T_\mathrm{ini},-1]}$ by $\bar y_{[-T_\mathrm{ini},-1]}+\sigma$ and adding $g_1(\sigma)=\lambda_\sigma\|\sigma\|_2^2$ and $g_2(\alpha)=\lambda_\alpha\|\alpha\|_2^2$ to the cost \cite{CoulsonLygDoe2019,BerberichKoeMueAll2021}. \autoref{fig:tutorial-deepc-bumpy} repeats the data-length experiment for $k=0.0025$ with data collected in closed loop around the target ($\lambda_\sigma=10^4$, $\lambda_\alpha$ scaled with the data energy, ten independent records per~$N$). The bound now only marks where the QP can be posed: at $N=51$ the median deviation is $0.7\,\mathrm{m}$ and individual records run away, the spread closes only around $N\approx150$, i.e.\ three times the bound, and the deviation then settles at $0.18\,\mathrm{m}$ instead of zero. This floor is structural: an \ac{LTI} predictor has a single equilibrium and a single constant input offset, so DeePC---like the \ac{MPC} built from the flat-terrain model---comes to rest one terrain valley beyond the target, a local minimum the true-model \ac{MPC} climbs out of.

  The regularization weights decide whether one even reaches that floor. \autoref{fig:tutorial-deepc-reg} sweeps $\lambda_\alpha$ over eight decades for three values of~$\lambda_\sigma$ on a terrain with $k=0.005$, using flat-terrain data. Only a tight history fit ($\lambda_\sigma=10^5$) with $\lambda_\alpha\le1$ recovers the flat-model behaviour, with a deviation of $0.20\,\mathrm{m}$ against $0.22\,\mathrm{m}$ for the \ac{MPC} with the flat model; with a looser fit ($\lambda_\sigma\le10^4$) the optimizer trades consistency against cost and, unless $\lambda_\alpha$ compensates, the car runs away down the terrain, while $\lambda_\alpha\ge10^2$ biases the prediction towards $\alpha=0$ and degrades the loop to $1$--$5\,\mathrm{m}$. Three conclusions follow. In the ideal \ac{LTI} case the data requirement is a counting bound and DeePC is equivalent to \ac{MPC}. Under noise or nonlinearity the bound is only necessary: regularization becomes indispensable and several times more data are needed before the result stabilizes. Finally, tuning matters as much as data: a well-tuned direct controller matches or slightly improves on an \ac{MPC} built from imperfect model knowledge---consistent with regularized DeePC being identification and control solved jointly rather than sequentially \cite{DorflerCoulMarTAC23}, an advantage that grows with measurement noise \cite[Ex.~9.8]{cfr-nootbook}---but it does not outperform the \ac{MPC} with the true model.

  \begin{tutorialfigure}
      \includegraphics[width=\linewidth]
          {examples/deepc/tutorial_deepc_flat.pdf}
      \captionof{figure}{
          Flat terrain: DeePC versus \ac{MPC} with the exact model as a function of the offline data length~$N$. Left: RMS deviation of the closed-loop position from the exact-model \ac{MPC}; below the bound (dashed) every QP is infeasible, from the bound on DeePC coincides with \ac{MPC}, and the shaded region is where persistency of excitation is guaranteed. Right: trajectories for the circled data lengths.
      }
      \label{fig:tutorial-deepc-flat}
  \end{tutorialfigure}

  \begin{tutorialfigure}
      \includegraphics[width=\linewidth]
          {examples/deepc/tutorial_deepc_bumpy.pdf}
      \captionof{figure}{
          Bumpy terrain with data recorded on the same terrain: median and 10--90\,\% band of the RMS deviation over ten data records per~$N$ (left) and trajectories for the circled data lengths (right). The bound no longer suffices, and the deviation settles at a floor set by the plant mismatch, one terrain valley beyond the target.
      }
      \label{fig:tutorial-deepc-bumpy}
  \end{tutorialfigure}

  \begin{tutorialfigure}
      \includegraphics[width=\linewidth]
          {examples/deepc/tutorial_deepc_regularization.pdf}
      \captionof{figure}{
          Effect of the regularization weights on bumpy terrain with flat-terrain data. (a)~RMS deviation from the exact-model \ac{MPC} as a function of $\lambda_\alpha$ for three values of $\lambda_\sigma$, against the \ac{MPC} with the flat-terrain model (dashed); (b)~trajectories for the circled settings. Too little regularization runs away, too much biases the prediction; the best setting matches the flat-model \ac{MPC} but not the exact-model one.
      }
      \label{fig:tutorial-deepc-reg}
  \end{tutorialfigure}
  \end{tutorialexample}
